
\documentclass{article}
\usepackage{colortbl}
\usepackage{makecell}
\usepackage{multirow}
\usepackage{supertabular}

\begin{document}

\newcounter{utterance}

\centering \large Interaction Transcript for game `clean\_up', experiment `1\_medium\_3obj\_en', episode 2 with nemotron{-}nano{-}9b{-}v2{-}no{-}reasoning{-}t1.0.
\vspace{24pt}

{ \footnotesize  \setcounter{utterance}{1}
\setlength{\tabcolsep}{0pt}
\begin{supertabular}{c@{$\;$}|p{.15\linewidth}@{}p{.15\linewidth}p{.15\linewidth}p{.15\linewidth}p{.15\linewidth}p{.15\linewidth}}
    \# & $\;$A & \multicolumn{4}{c}{Game Master} & $\;\:$B\\
    \hline

    \theutterance \stepcounter{utterance}  
    & & \multicolumn{4}{p{0.6\linewidth}}{
        \cellcolor[rgb]{0.9,0.9,0.9}{
            \makecell[{{p{\linewidth}}}]{
                \texttt{\tiny{[P1$\langle$GM]}}
                \texttt{Sono il tuo game master, e stai giocando un gioco collaborativo con la seguente griglia come tavolo da gioco:} \\
\\ 
\texttt{```} \\
\texttt{ 1234567} \\
\texttt{╔═╤═╤═╤═╗} \\
\texttt{║◌│◌│◌│◌║ 1} \\
\texttt{╟─┴─┼─┤◌║ 2} \\
\texttt{║C◌◌│◌└─╢ 3} \\
\texttt{╟───┘◌◌◌║ 4} \\
\texttt{║◌L◌◌◌┌─╢ 5} \\
\texttt{║◌◌◌◌◌│◌║ 6} \\
\texttt{║◌◌◌◌P│◌║ 7} \\
\texttt{╚═════╧═╝} \\
\\ 
\texttt{```} \\
\\ 
\texttt{* Il bordo superiore mostra le coordinate x che aumentano verso destra, e il bordo destro le coordinate y che aumentano verso il basso.} \\
\texttt{* I seguenti oggetti sono posizionati casualmente sulla tua griglia: 'C', 'L', 'P'.} \\
\\ 
\\ 
\texttt{L'altro giocatore vede una variante del tavolo da gioco, dove gli oggetti sono posizionati in diverse posizioni casuali. Non puoi vedere il tavolo dell'altro giocatore, e loro non possono vedere il tuo.} \\
\\ 
\\ 
\texttt{**Obiettivo:**} \\
\\ 
\texttt{Entrambi i giocatori devono spostare gli oggetti sui rispettivi sfondi in modo che oggetti identici finiscano alle stesse coordinate. Devi comunicare con l'altro giocatore per concordare una configurazione obiettivo comune.} \\
\\ 
\\ 
\texttt{**Regole:**} \\
\\ 
\texttt{* In ogni turno, puoi inviare esattamente uno dei seguenti due comandi:} \\
\texttt{ 1. `COMUNICA: <MESSAGE>`: per inviare un messaggio (tutto fino alla prossima interruzione di riga) all'altro giocatore. Lo inoltrerò al tuo partner.} \\
\texttt{2. `SPOSTA: <OBJECT>, (<X>, <Y>)`: per spostare un oggetto in una nuova posizione, dove <X> è la colonna e <Y> è la riga. Ti farò sapere se la tua mossa è andata a buon fine o no.} \\
\texttt{* Se non rispetti il formato, o invii diversi comandi contemporaneamente, devo penalizzarti.} \\
\texttt{* Se entrambi i giocatori accumulano più di 8 penalità, perdete entrambi il gioco.} \\
\texttt{* È vitale che comunichi con l'altro giocatore riguardo al tuo stato obiettivo! L'*unico* modo in cui puoi trasmettere la tua strategia all'altro giocatore è usando il comando `COMUNICA: <MESSAGE>`!} \\
\\ 
\\ 
\texttt{**Spostamento Oggetti**} \\
\\ 
\texttt{* Puoi spostare oggetti solo in celle entro i limiti della griglia. La cella di destinazione deve essere vuota, cioè deve contenere solo il simbolo '◌'.} \\
\texttt{* Se provi a spostare un oggetto in un posto che non è vuoto, o provi a spostarlo fuori dalla griglia, devo penalizzarti. Hai un altro tentativo.} \\
\texttt{* Prima di fare una mossa, controlla due volte che il posto di destinazione sia vuoto, e non contenga alcuna lettera, cornice, o linea!} \\
\\ 
\\ 
\texttt{**Fine del Gioco**} \\
\\ 
\texttt{Se pensi di aver raggiunto l'obiettivo di allineare tutti gli oggetti, puoi chiedere all'altro giocatore di finire il gioco inviando COMUNICA: finito?. Se l'altro giocatore ti chiede di finire il gioco, e rispondi con COMUNICA: finito!, the game will end.} \\
\\ 
\texttt{Entrambi i giocatori vincono se il gioco finisce entro 12 round, dove un round è definito come due giocatori che inviano ciascuno un comando valido.} \\
\\ 
\texttt{**Punteggio:**} \\
\\ 
\texttt{Più vicini sono gli oggetti identici in entrambi i tavoli da gioco, più punti ottieni. Le penalità riducono i tuoi punti. Riesci a battere il record?} \\
\\ 
\\ 
\texttt{Per favore invia un messaggio all'altro giocatore per iniziare il gioco!} \\
            }
        }
    }
    & \\ \\

    \theutterance \stepcounter{utterance}  
    & \multicolumn{4}{p{0.6\linewidth}}{
        \cellcolor[rgb]{0.9,0.9,0.9}{
            \makecell[{{p{\linewidth}}}]{
                \texttt{\tiny{[P1$\rangle$GM]}}
                \texttt{COMUNICA: Ciao! Sto cercando di coordinare le posizioni di C, L e P. Possiamo concordare un piano per spostare gli oggetti verso le posizioni centrali o un angolo specifico? Vi proponi un punto di riferimento? 😊} \\
            }
        }
    }
    & & \\ \\

    \theutterance \stepcounter{utterance}  
    & & & \multicolumn{2}{p{0.3\linewidth}}{
        \cellcolor[rgb]{0.9,0.9,0.9}{
            \makecell[{{p{\linewidth}}}]{
                \texttt{\tiny{[GM$|$GM]}}
                \texttt{COMUNICA: Ciao! Sto cercando di coordinare le posizioni di C, L e P. Possiamo concordare un piano per spostare gli oggetti verso le posizioni centrali o un angolo specifico? Vi proponi un punto di riferimento? 😊} \\
            }
        }
    }
    & & \\ \\

    \theutterance \stepcounter{utterance}  
    & & \multicolumn{4}{p{0.6\linewidth}}{
        \cellcolor[rgb]{0.9,0.9,0.9}{
            \makecell[{{p{\linewidth}}}]{
                \texttt{\tiny{[GM$\rangle$P2]}}
                \texttt{Sono il tuo game master, e stai giocando un gioco collaborativo con la seguente griglia come tavolo da gioco:} \\
\\ 
\texttt{```} \\
\texttt{ 1234567} \\
\texttt{╔═╤═╤═╤═╗} \\
\texttt{║◌│◌│L│◌║ 1} \\
\texttt{╟─┴─┼─┤◌║ 2} \\
\texttt{║C◌◌│◌└─╢ 3} \\
\texttt{╟───┘◌◌◌║ 4} \\
\texttt{║P◌◌◌◌┌─╢ 5} \\
\texttt{║◌◌◌◌◌│◌║ 6} \\
\texttt{║◌◌◌◌◌│◌║ 7} \\
\texttt{╚═════╧═╝} \\
\\ 
\texttt{```} \\
\\ 
\texttt{* Il bordo superiore mostra le coordinate x che aumentano verso destra, e il bordo destro le coordinate y che aumentano verso il basso.} \\
\texttt{* I seguenti oggetti sono posizionati casualmente sulla tua griglia: 'C', 'L', 'P'.} \\
\\ 
\\ 
\texttt{L'altro giocatore vede una variante del tavolo da gioco, dove gli oggetti sono posizionati in diverse posizioni casuali. Non puoi vedere il tavolo dell'altro giocatore, e loro non possono vedere il tuo.} \\
\\ 
\\ 
\texttt{**Obiettivo:**} \\
\\ 
\texttt{Entrambi i giocatori devono spostare gli oggetti sui rispettivi sfondi in modo che oggetti identici finiscano alle stesse coordinate. Devi comunicare con l'altro giocatore per concordare una configurazione obiettivo comune.} \\
\\ 
\\ 
\texttt{**Regole:**} \\
\\ 
\texttt{* In ogni turno, puoi inviare esattamente uno dei seguenti due comandi:} \\
\texttt{ 1. `COMUNICA: <MESSAGE>`: per inviare un messaggio (tutto fino alla prossima interruzione di riga) all'altro giocatore. Lo inoltrerò al tuo partner.} \\
\texttt{2. `SPOSTA: <OBJECT>, (<X>, <Y>)`: per spostare un oggetto in una nuova posizione, dove <X> è la colonna e <Y> è la riga. Ti farò sapere se la tua mossa è andata a buon fine o no.} \\
\texttt{* Se non rispetti il formato, o invii diversi comandi contemporaneamente, devo penalizzarti.} \\
\texttt{* Se entrambi i giocatori accumulano più di 8 penalità, perdete entrambi il gioco.} \\
\texttt{* È vitale che comunichi con l'altro giocatore riguardo al tuo stato obiettivo! L'*unico* modo in cui puoi trasmettere la tua strategia all'altro giocatore è usando il comando `COMUNICA: <MESSAGE>`!} \\
\\ 
\\ 
\texttt{**Spostamento Oggetti**} \\
\\ 
\texttt{* Puoi spostare oggetti solo in celle entro i limiti della griglia. La cella di destinazione deve essere vuota, cioè deve contenere solo il simbolo '◌'.} \\
\texttt{* Se provi a spostare un oggetto in un posto che non è vuoto, o provi a spostarlo fuori dalla griglia, devo penalizzarti. Hai un altro tentativo.} \\
\texttt{* Prima di fare una mossa, controlla due volte che il posto di destinazione sia vuoto, e non contenga alcuna lettera, cornice, o linea!} \\
\\ 
\\ 
\texttt{**Fine del Gioco**} \\
\\ 
\texttt{Se pensi di aver raggiunto l'obiettivo di allineare tutti gli oggetti, puoi chiedere all'altro giocatore di finire il gioco inviando COMUNICA: finito?. Se l'altro giocatore ti chiede di finire il gioco, e rispondi con COMUNICA: finito!, the game will end.} \\
\\ 
\texttt{Entrambi i giocatori vincono se il gioco finisce entro 12 round, dove un round è definito come due giocatori che inviano ciascuno un comando valido.} \\
\\ 
\texttt{**Punteggio:**} \\
\\ 
\texttt{Più vicini sono gli oggetti identici in entrambi i tavoli da gioco, più punti ottieni. Le penalità riducono i tuoi punti. Riesci a battere il record?} \\
\\ 
\texttt{L'altro giocatore ha iniziato il gioco inviando questo messaggio:} \\
\\ 
\texttt{"Ciao! Sto cercando di coordinare le posizioni di C, L e P. Possiamo concordare un piano per spostare gli oggetti verso le posizioni centrali o un angolo specifico? Vi proponi un punto di riferimento? 😊"} \\
\\ 
\texttt{Qual è il tuo primo comando?} \\
            }
        }
    }
    & \\ \\

    \theutterance \stepcounter{utterance}  
    & & & \multicolumn{4}{p{0.6\linewidth}}{
        \cellcolor[rgb]{0.9,0.9,0.9}{
            \makecell[{{p{\linewidth}}}]{
                \texttt{\tiny{[GM$\langle$P2]}}
                \texttt{COMUNICA: Ciao! Proponiamo di spostare tutti gli oggetti (C, L, P) al centro della griglia, specificamente alla posizione (4,4). Questo punto è equidistante e facile da raggiungere per entrambi. Possiamo poi coordinare i movimenti passo dopo passo. 😊} \\
            }
        }
    }
     \\ \\

    \theutterance \stepcounter{utterance}  
    & & & \multicolumn{2}{p{0.3\linewidth}}{
        \cellcolor[rgb]{0.9,0.9,0.9}{
            \makecell[{{p{\linewidth}}}]{
                \texttt{\tiny{[GM$|$GM]}}
                \texttt{COMUNICA: Ciao! Proponiamo di spostare tutti gli oggetti (C, L, P) al centro della griglia, specificamente alla posizione (4,4). Questo punto è equidistante e facile da raggiungere per entrambi. Possiamo poi coordinare i movimenti passo dopo passo. 😊} \\
            }
        }
    }
    & & \\ \\

    \theutterance \stepcounter{utterance}  
    & & \multicolumn{4}{p{0.6\linewidth}}{
        \cellcolor[rgb]{0.9,0.9,0.9}{
            \makecell[{{p{\linewidth}}}]{
                \texttt{\tiny{[P1$\langle$GM]}}
                \texttt{Il tuo messaggio è stato inoltrato all'altro giocatore.} \\
\\ 
\\ 
\texttt{Stai attualmente giocando il round 1 di massimo 12.} \\
\texttt{Avete accumulato collettivamente 0 di 8 penalità.} \\
\texttt{L'altro giocatore ha inviato questo messaggio:} \\
\\ 
\\ 
\texttt{"Ciao! Proponiamo di spostare tutti gli oggetti (C, L, P) al centro della griglia, specificamente alla posizione (4,4). Questo punto è equidistante e facile da raggiungere per entrambi. Possiamo poi coordinare i movimenti passo dopo passo. 😊"} \\
\\ 
\\ 
\texttt{Qual è il tuo prossimo comando?} \\
            }
        }
    }
    & \\ \\

    \theutterance \stepcounter{utterance}  
    & \multicolumn{4}{p{0.6\linewidth}}{
        \cellcolor[rgb]{0.9,0.9,0.9}{
            \makecell[{{p{\linewidth}}}]{
                \texttt{\tiny{[P1$\rangle$GM]}}
                \texttt{COMUNICA: Accordo con il piano di spostare tutti gli oggetti al centro (4,4). Possiamo iniziare con 'C' o 'L'? Vi proponi di muovere prima 'C' verso (4,4) o 'L'? 😊} \\
            }
        }
    }
    & & \\ \\

    \theutterance \stepcounter{utterance}  
    & & & \multicolumn{2}{p{0.3\linewidth}}{
        \cellcolor[rgb]{0.9,0.9,0.9}{
            \makecell[{{p{\linewidth}}}]{
                \texttt{\tiny{[GM$|$GM]}}
                \texttt{COMUNICA: Accordo con il piano di spostare tutti gli oggetti al centro (4,4). Possiamo iniziare con 'C' o 'L'? Vi proponi di muovere prima 'C' verso (4,4) o 'L'? 😊} \\
            }
        }
    }
    & & \\ \\

    \theutterance \stepcounter{utterance}  
    & & \multicolumn{4}{p{0.6\linewidth}}{
        \cellcolor[rgb]{0.9,0.9,0.9}{
            \makecell[{{p{\linewidth}}}]{
                \texttt{\tiny{[GM$\rangle$P2]}}
                \texttt{Il tuo messaggio è stato inoltrato all'altro giocatore.} \\
\\ 
\\ 
\texttt{Stai attualmente giocando il round 2 di massimo 12.} \\
\texttt{Avete accumulato collettivamente 0 di 8 penalità.} \\
\texttt{L'altro giocatore ha inviato questo messaggio:} \\
\\ 
\\ 
\texttt{"Accordo con il piano di spostare tutti gli oggetti al centro (4,4). Possiamo iniziare con 'C' o 'L'? Vi proponi di muovere prima 'C' verso (4,4) o 'L'? 😊"} \\
\\ 
\\ 
\texttt{Qual è il tuo prossimo comando?} \\
            }
        }
    }
    & \\ \\

    \theutterance \stepcounter{utterance}  
    & & & \multicolumn{4}{p{0.6\linewidth}}{
        \cellcolor[rgb]{0.9,0.9,0.9}{
            \makecell[{{p{\linewidth}}}]{
                \texttt{\tiny{[GM$\langle$P2]}}
                \texttt{COMUNICA: Proponiamo di muovere prima 'C' verso (4,4). Se la posizione è occupata, possiamo adattare il piano. 😊} \\
            }
        }
    }
     \\ \\

    \theutterance \stepcounter{utterance}  
    & & & \multicolumn{2}{p{0.3\linewidth}}{
        \cellcolor[rgb]{0.9,0.9,0.9}{
            \makecell[{{p{\linewidth}}}]{
                \texttt{\tiny{[GM$|$GM]}}
                \texttt{COMUNICA: Proponiamo di muovere prima 'C' verso (4,4). Se la posizione è occupata, possiamo adattare il piano. 😊} \\
            }
        }
    }
    & & \\ \\

    \theutterance \stepcounter{utterance}  
    & & \multicolumn{4}{p{0.6\linewidth}}{
        \cellcolor[rgb]{0.9,0.9,0.9}{
            \makecell[{{p{\linewidth}}}]{
                \texttt{\tiny{[P1$\langle$GM]}}
                \texttt{Il tuo messaggio è stato inoltrato all'altro giocatore.} \\
\\ 
\\ 
\texttt{Stai attualmente giocando il round 2 di massimo 12.} \\
\texttt{Avete accumulato collettivamente 0 di 8 penalità.} \\
\texttt{L'altro giocatore ha inviato questo messaggio:} \\
\\ 
\\ 
\texttt{"Proponiamo di muovere prima 'C' verso (4,4). Se la posizione è occupata, possiamo adattare il piano. 😊"} \\
\\ 
\\ 
\texttt{Qual è il tuo prossimo comando?} \\
            }
        }
    }
    & \\ \\

    \theutterance \stepcounter{utterance}  
    & \multicolumn{4}{p{0.6\linewidth}}{
        \cellcolor[rgb]{0.9,0.9,0.9}{
            \makecell[{{p{\linewidth}}}]{
                \texttt{\tiny{[P1$\rangle$GM]}}
                \texttt{SPOSTA: C, (4,4)} \\
            }
        }
    }
    & & \\ \\

    \theutterance \stepcounter{utterance}  
    & & & \multicolumn{2}{p{0.3\linewidth}}{
        \cellcolor[rgb]{0.9,0.9,0.9}{
            \makecell[{{p{\linewidth}}}]{
                \texttt{\tiny{[GM$|$GM]}}
                \texttt{SPOSTA: C, (4,4)} \\
            }
        }
    }
    & & \\ \\

    \theutterance \stepcounter{utterance}  
    & & & \multicolumn{2}{p{0.3\linewidth}}{
        \cellcolor[rgb]{0.9,0.9,0.9}{
            \makecell[{{p{\linewidth}}}]{
                \texttt{\tiny{[GM$|$GM]}}
                \texttt{Penalità: (4,4) non è vuoto, ma contiene '┘'.} \\
\\ 
\texttt{Avete accumulato collettivamente 1 di 8 penalità.} \\
\texttt{Riprova!} \\
            }
        }
    }
    & & \\ \\

    \theutterance \stepcounter{utterance}  
    & & \multicolumn{4}{p{0.6\linewidth}}{
        \cellcolor[rgb]{0.9,0.9,0.9}{
            \makecell[{{p{\linewidth}}}]{
                \texttt{\tiny{[P1$\langle$GM]}}
                \texttt{Penalità: (4,4) non è vuoto, ma contiene '┘'.} \\
\\ 
\texttt{Avete accumulato collettivamente 1 di 8 penalità.} \\
\texttt{Riprova!} \\
            }
        }
    }
    & \\ \\

    \theutterance \stepcounter{utterance}  
    & \multicolumn{4}{p{0.6\linewidth}}{
        \cellcolor[rgb]{0.9,0.9,0.9}{
            \makecell[{{p{\linewidth}}}]{
                \texttt{\tiny{[P1$\rangle$GM]}}
                \texttt{COMUNICA: Scusa per l'errore! La posizione (4,4) è occupata. Proponiamo di spostare 'C' in una cella vicina e libera, come (3,3) o (5,3). Quale posizione preferisci? 😊} \\
            }
        }
    }
    & & \\ \\

    \theutterance \stepcounter{utterance}  
    & & & \multicolumn{2}{p{0.3\linewidth}}{
        \cellcolor[rgb]{0.9,0.9,0.9}{
            \makecell[{{p{\linewidth}}}]{
                \texttt{\tiny{[GM$|$GM]}}
                \texttt{COMUNICA: Scusa per l'errore! La posizione (4,4) è occupata. Proponiamo di spostare 'C' in una cella vicina e libera, come (3,3) o (5,3). Quale posizione preferisci? 😊} \\
            }
        }
    }
    & & \\ \\

    \theutterance \stepcounter{utterance}  
    & & \multicolumn{4}{p{0.6\linewidth}}{
        \cellcolor[rgb]{0.9,0.9,0.9}{
            \makecell[{{p{\linewidth}}}]{
                \texttt{\tiny{[GM$\rangle$P2]}}
                \texttt{Il tuo messaggio è stato inoltrato all'altro giocatore.} \\
\\ 
\\ 
\texttt{Stai attualmente giocando il round 3 di massimo 12.} \\
\texttt{Avete accumulato collettivamente 1 di 8 penalità.} \\
\texttt{L'altro giocatore ha inviato questo messaggio:} \\
\\ 
\\ 
\texttt{"Scusa per l'errore! La posizione (4,4) è occupata. Proponiamo di spostare 'C' in una cella vicina e libera, come (3,3) o (5,3). Quale posizione preferisci? 😊"} \\
\\ 
\\ 
\texttt{Qual è il tuo prossimo comando?} \\
            }
        }
    }
    & \\ \\

    \theutterance \stepcounter{utterance}  
    & & & \multicolumn{4}{p{0.6\linewidth}}{
        \cellcolor[rgb]{0.9,0.9,0.9}{
            \makecell[{{p{\linewidth}}}]{
                \texttt{\tiny{[GM$\langle$P2]}}
                \texttt{COMUNICA: Accordo di spostare 'C' a (3,3) per evitare conflitti. Procediamo con questa posizione. 😊} \\
            }
        }
    }
     \\ \\

    \theutterance \stepcounter{utterance}  
    & & & \multicolumn{2}{p{0.3\linewidth}}{
        \cellcolor[rgb]{0.9,0.9,0.9}{
            \makecell[{{p{\linewidth}}}]{
                \texttt{\tiny{[GM$|$GM]}}
                \texttt{COMUNICA: Accordo di spostare 'C' a (3,3) per evitare conflitti. Procediamo con questa posizione. 😊} \\
            }
        }
    }
    & & \\ \\

    \theutterance \stepcounter{utterance}  
    & & \multicolumn{4}{p{0.6\linewidth}}{
        \cellcolor[rgb]{0.9,0.9,0.9}{
            \makecell[{{p{\linewidth}}}]{
                \texttt{\tiny{[P1$\langle$GM]}}
                \texttt{Il tuo messaggio è stato inoltrato all'altro giocatore.} \\
\\ 
\\ 
\texttt{Stai attualmente giocando il round 3 di massimo 12.} \\
\texttt{Avete accumulato collettivamente 1 di 8 penalità.} \\
\texttt{L'altro giocatore ha inviato questo messaggio:} \\
\\ 
\\ 
\texttt{"Accordo di spostare 'C' a (3,3) per evitare conflitti. Procediamo con questa posizione. 😊"} \\
\\ 
\\ 
\texttt{Qual è il tuo prossimo comando?} \\
            }
        }
    }
    & \\ \\

    \theutterance \stepcounter{utterance}  
    & \multicolumn{4}{p{0.6\linewidth}}{
        \cellcolor[rgb]{0.9,0.9,0.9}{
            \makecell[{{p{\linewidth}}}]{
                \texttt{\tiny{[P1$\rangle$GM]}}
                \texttt{SPOSTA: C, (3,3)} \\
            }
        }
    }
    & & \\ \\

    \theutterance \stepcounter{utterance}  
    & & & \multicolumn{2}{p{0.3\linewidth}}{
        \cellcolor[rgb]{0.9,0.9,0.9}{
            \makecell[{{p{\linewidth}}}]{
                \texttt{\tiny{[GM$|$GM]}}
                \texttt{SPOSTA: C, (3,3)} \\
            }
        }
    }
    & & \\ \\

    \theutterance \stepcounter{utterance}  
    & & & \multicolumn{2}{p{0.3\linewidth}}{
        \cellcolor[rgb]{0.9,0.9,0.9}{
            \makecell[{{p{\linewidth}}}]{
                \texttt{\tiny{[GM$|$GM]}}
                \texttt{Spostato 'C' in (3,3) con successo. La tua griglia aggiornata appare così:} \\
\texttt{```} \\
\texttt{ 1234567} \\
\texttt{╔═╤═╤═╤═╗} \\
\texttt{║◌│◌│◌│◌║ 1} \\
\texttt{╟─┴─┼─┤◌║ 2} \\
\texttt{║◌◌C│◌└─╢ 3} \\
\texttt{╟───┘◌◌◌║ 4} \\
\texttt{║◌L◌◌◌┌─╢ 5} \\
\texttt{║◌◌◌◌◌│◌║ 6} \\
\texttt{║◌◌◌◌P│◌║ 7} \\
\texttt{╚═════╧═╝} \\
\texttt{```} \\
            }
        }
    }
    & & \\ \\

    \theutterance \stepcounter{utterance}  
    & & \multicolumn{4}{p{0.6\linewidth}}{
        \cellcolor[rgb]{0.9,0.9,0.9}{
            \makecell[{{p{\linewidth}}}]{
                \texttt{\tiny{[GM$\rangle$P2]}}
                \texttt{Il tuo messaggio è stato inoltrato all'altro giocatore.} \\
\\ 
\\ 
\texttt{Stai attualmente giocando il round 4 di massimo 12.} \\
\texttt{Avete accumulato collettivamente 1 di 8 penalità.} \\
\texttt{L'altro giocatore ha spostato un oggetto sulla sua griglia.} \\
\\ 
\\ 
\texttt{Qual è il tuo prossimo comando?} \\
            }
        }
    }
    & \\ \\

    \theutterance \stepcounter{utterance}  
    & & & \multicolumn{4}{p{0.6\linewidth}}{
        \cellcolor[rgb]{0.9,0.9,0.9}{
            \makecell[{{p{\linewidth}}}]{
                \texttt{\tiny{[GM$\langle$P2]}}
                \texttt{SPOSTA: C, (3,3)} \\
            }
        }
    }
     \\ \\

    \theutterance \stepcounter{utterance}  
    & & & \multicolumn{2}{p{0.3\linewidth}}{
        \cellcolor[rgb]{0.9,0.9,0.9}{
            \makecell[{{p{\linewidth}}}]{
                \texttt{\tiny{[GM$|$GM]}}
                \texttt{SPOSTA: C, (3,3)} \\
            }
        }
    }
    & & \\ \\

    \theutterance \stepcounter{utterance}  
    & & & \multicolumn{2}{p{0.3\linewidth}}{
        \cellcolor[rgb]{0.9,0.9,0.9}{
            \makecell[{{p{\linewidth}}}]{
                \texttt{\tiny{[GM$|$GM]}}
                \texttt{Spostato 'C' in (3,3) con successo. La tua griglia aggiornata appare così:} \\
\texttt{```} \\
\texttt{ 1234567} \\
\texttt{╔═╤═╤═╤═╗} \\
\texttt{║◌│◌│L│◌║ 1} \\
\texttt{╟─┴─┼─┤◌║ 2} \\
\texttt{║◌◌C│◌└─╢ 3} \\
\texttt{╟───┘◌◌◌║ 4} \\
\texttt{║P◌◌◌◌┌─╢ 5} \\
\texttt{║◌◌◌◌◌│◌║ 6} \\
\texttt{║◌◌◌◌◌│◌║ 7} \\
\texttt{╚═════╧═╝} \\
\texttt{```} \\
            }
        }
    }
    & & \\ \\

    \theutterance \stepcounter{utterance}  
    & & \multicolumn{4}{p{0.6\linewidth}}{
        \cellcolor[rgb]{0.9,0.9,0.9}{
            \makecell[{{p{\linewidth}}}]{
                \texttt{\tiny{[P1$\langle$GM]}}
                \texttt{Spostato 'C' in (3,3) con successo. La tua griglia aggiornata appare così:} \\
\texttt{```} \\
\texttt{ 1234567} \\
\texttt{╔═╤═╤═╤═╗} \\
\texttt{║◌│◌│◌│◌║ 1} \\
\texttt{╟─┴─┼─┤◌║ 2} \\
\texttt{║◌◌C│◌└─╢ 3} \\
\texttt{╟───┘◌◌◌║ 4} \\
\texttt{║◌L◌◌◌┌─╢ 5} \\
\texttt{║◌◌◌◌◌│◌║ 6} \\
\texttt{║◌◌◌◌P│◌║ 7} \\
\texttt{╚═════╧═╝} \\
\texttt{```} \\
\texttt{Stai attualmente giocando il round 4 di massimo 12.} \\
\texttt{Avete accumulato collettivamente 1 di 8 penalità.} \\
\texttt{L'altro giocatore ha spostato un oggetto sulla sua griglia.} \\
\\ 
\\ 
\texttt{Qual è il tuo prossimo comando?} \\
            }
        }
    }
    & \\ \\

    \theutterance \stepcounter{utterance}  
    & \multicolumn{4}{p{0.6\linewidth}}{
        \cellcolor[rgb]{0.9,0.9,0.9}{
            \makecell[{{p{\linewidth}}}]{
                \texttt{\tiny{[P1$\rangle$GM]}}
                \texttt{COMUNICA: Ottimo! Ora che 'C' è a (3,3), possiamo coordinare i prossimi passi. Proponiamo di spostare 'L' verso (4,4) o 'P' verso (4,4) nel turno successivo. Quale oggetto vuoi muovere per primo? 😊} \\
            }
        }
    }
    & & \\ \\

    \theutterance \stepcounter{utterance}  
    & & & \multicolumn{2}{p{0.3\linewidth}}{
        \cellcolor[rgb]{0.9,0.9,0.9}{
            \makecell[{{p{\linewidth}}}]{
                \texttt{\tiny{[GM$|$GM]}}
                \texttt{COMUNICA: Ottimo! Ora che 'C' è a (3,3), possiamo coordinare i prossimi passi. Proponiamo di spostare 'L' verso (4,4) o 'P' verso (4,4) nel turno successivo. Quale oggetto vuoi muovere per primo? 😊} \\
            }
        }
    }
    & & \\ \\

    \theutterance \stepcounter{utterance}  
    & & \multicolumn{4}{p{0.6\linewidth}}{
        \cellcolor[rgb]{0.9,0.9,0.9}{
            \makecell[{{p{\linewidth}}}]{
                \texttt{\tiny{[GM$\rangle$P2]}}
                \texttt{Spostato 'C' in (3,3) con successo. La tua griglia aggiornata appare così:} \\
\texttt{```} \\
\texttt{ 1234567} \\
\texttt{╔═╤═╤═╤═╗} \\
\texttt{║◌│◌│L│◌║ 1} \\
\texttt{╟─┴─┼─┤◌║ 2} \\
\texttt{║◌◌C│◌└─╢ 3} \\
\texttt{╟───┘◌◌◌║ 4} \\
\texttt{║P◌◌◌◌┌─╢ 5} \\
\texttt{║◌◌◌◌◌│◌║ 6} \\
\texttt{║◌◌◌◌◌│◌║ 7} \\
\texttt{╚═════╧═╝} \\
\texttt{```} \\
\texttt{Stai attualmente giocando il round 5 di massimo 12.} \\
\texttt{Avete accumulato collettivamente 1 di 8 penalità.} \\
\texttt{L'altro giocatore ha inviato questo messaggio:} \\
\\ 
\\ 
\texttt{"Ottimo! Ora che 'C' è a (3,3), possiamo coordinare i prossimi passi. Proponiamo di spostare 'L' verso (4,4) o 'P' verso (4,4) nel turno successivo. Quale oggetto vuoi muovere per primo? 😊"} \\
\\ 
\\ 
\texttt{Qual è il tuo prossimo comando?} \\
            }
        }
    }
    & \\ \\

    \theutterance \stepcounter{utterance}  
    & & & \multicolumn{4}{p{0.6\linewidth}}{
        \cellcolor[rgb]{0.9,0.9,0.9}{
            \makecell[{{p{\linewidth}}}]{
                \texttt{\tiny{[GM$\langle$P2]}}
                \texttt{COMUNICA: Proponiamo di spostare 'L' verso (4,4) per primo. Se la posizione è libera, procediamo; altrimenti, adattiamo il piano. 😊} \\
            }
        }
    }
     \\ \\

    \theutterance \stepcounter{utterance}  
    & & & \multicolumn{2}{p{0.3\linewidth}}{
        \cellcolor[rgb]{0.9,0.9,0.9}{
            \makecell[{{p{\linewidth}}}]{
                \texttt{\tiny{[GM$|$GM]}}
                \texttt{COMUNICA: Proponiamo di spostare 'L' verso (4,4) per primo. Se la posizione è libera, procediamo; altrimenti, adattiamo il piano. 😊} \\
            }
        }
    }
    & & \\ \\

    \theutterance \stepcounter{utterance}  
    & & \multicolumn{4}{p{0.6\linewidth}}{
        \cellcolor[rgb]{0.9,0.9,0.9}{
            \makecell[{{p{\linewidth}}}]{
                \texttt{\tiny{[P1$\langle$GM]}}
                \texttt{Il tuo messaggio è stato inoltrato all'altro giocatore.} \\
\\ 
\\ 
\texttt{Stai attualmente giocando il round 5 di massimo 12.} \\
\texttt{Avete accumulato collettivamente 1 di 8 penalità.} \\
\texttt{L'altro giocatore ha inviato questo messaggio:} \\
\\ 
\\ 
\texttt{"Proponiamo di spostare 'L' verso (4,4) per primo. Se la posizione è libera, procediamo; altrimenti, adattiamo il piano. 😊"} \\
\\ 
\\ 
\texttt{Qual è il tuo prossimo comando?} \\
            }
        }
    }
    & \\ \\

    \theutterance \stepcounter{utterance}  
    & \multicolumn{4}{p{0.6\linewidth}}{
        \cellcolor[rgb]{0.9,0.9,0.9}{
            \makecell[{{p{\linewidth}}}]{
                \texttt{\tiny{[P1$\rangle$GM]}}
                \texttt{SPOSTA: L, (4,4)} \\
            }
        }
    }
    & & \\ \\

    \theutterance \stepcounter{utterance}  
    & & & \multicolumn{2}{p{0.3\linewidth}}{
        \cellcolor[rgb]{0.9,0.9,0.9}{
            \makecell[{{p{\linewidth}}}]{
                \texttt{\tiny{[GM$|$GM]}}
                \texttt{SPOSTA: L, (4,4)} \\
            }
        }
    }
    & & \\ \\

    \theutterance \stepcounter{utterance}  
    & & & \multicolumn{2}{p{0.3\linewidth}}{
        \cellcolor[rgb]{0.9,0.9,0.9}{
            \makecell[{{p{\linewidth}}}]{
                \texttt{\tiny{[GM$|$GM]}}
                \texttt{Penalità: (4,4) non è vuoto, ma contiene '┘'.} \\
\\ 
\texttt{Avete accumulato collettivamente 2 di 8 penalità.} \\
\texttt{Riprova!} \\
            }
        }
    }
    & & \\ \\

    \theutterance \stepcounter{utterance}  
    & & \multicolumn{4}{p{0.6\linewidth}}{
        \cellcolor[rgb]{0.9,0.9,0.9}{
            \makecell[{{p{\linewidth}}}]{
                \texttt{\tiny{[P1$\langle$GM]}}
                \texttt{Penalità: (4,4) non è vuoto, ma contiene '┘'.} \\
\\ 
\texttt{Avete accumulato collettivamente 2 di 8 penalità.} \\
\texttt{Riprova!} \\
            }
        }
    }
    & \\ \\

    \theutterance \stepcounter{utterance}  
    & \multicolumn{4}{p{0.6\linewidth}}{
        \cellcolor[rgb]{0.9,0.9,0.9}{
            \makecell[{{p{\linewidth}}}]{
                \texttt{\tiny{[P1$\rangle$GM]}}
                \texttt{COMUNICA: Scusa per il problema! La posizione (4,4) è ancora occupata. Proponiamo di spostare 'L' in una cella vicina e libera, come (4,3) o (5,4). Quale posizione preferisci? 😊} \\
            }
        }
    }
    & & \\ \\

    \theutterance \stepcounter{utterance}  
    & & & \multicolumn{2}{p{0.3\linewidth}}{
        \cellcolor[rgb]{0.9,0.9,0.9}{
            \makecell[{{p{\linewidth}}}]{
                \texttt{\tiny{[GM$|$GM]}}
                \texttt{COMUNICA: Scusa per il problema! La posizione (4,4) è ancora occupata. Proponiamo di spostare 'L' in una cella vicina e libera, come (4,3) o (5,4). Quale posizione preferisci? 😊} \\
            }
        }
    }
    & & \\ \\

    \theutterance \stepcounter{utterance}  
    & & \multicolumn{4}{p{0.6\linewidth}}{
        \cellcolor[rgb]{0.9,0.9,0.9}{
            \makecell[{{p{\linewidth}}}]{
                \texttt{\tiny{[GM$\rangle$P2]}}
                \texttt{Il tuo messaggio è stato inoltrato all'altro giocatore.} \\
\\ 
\\ 
\texttt{Stai attualmente giocando il round 6 di massimo 12.} \\
\texttt{Avete accumulato collettivamente 2 di 8 penalità.} \\
\texttt{L'altro giocatore ha inviato questo messaggio:} \\
\\ 
\\ 
\texttt{"Scusa per il problema! La posizione (4,4) è ancora occupata. Proponiamo di spostare 'L' in una cella vicina e libera, come (4,3) o (5,4). Quale posizione preferisci? 😊"} \\
\\ 
\\ 
\texttt{Qual è il tuo prossimo comando?} \\
            }
        }
    }
    & \\ \\

    \theutterance \stepcounter{utterance}  
    & & & \multicolumn{4}{p{0.6\linewidth}}{
        \cellcolor[rgb]{0.9,0.9,0.9}{
            \makecell[{{p{\linewidth}}}]{
                \texttt{\tiny{[GM$\langle$P2]}}
                \texttt{COMUNICA: Accordo di spostare 'L' a (4,3) per evitare conflitti. Procediamo con questa posizione. 😊} \\
            }
        }
    }
     \\ \\

    \theutterance \stepcounter{utterance}  
    & & & \multicolumn{2}{p{0.3\linewidth}}{
        \cellcolor[rgb]{0.9,0.9,0.9}{
            \makecell[{{p{\linewidth}}}]{
                \texttt{\tiny{[GM$|$GM]}}
                \texttt{COMUNICA: Accordo di spostare 'L' a (4,3) per evitare conflitti. Procediamo con questa posizione. 😊} \\
            }
        }
    }
    & & \\ \\

    \theutterance \stepcounter{utterance}  
    & & \multicolumn{4}{p{0.6\linewidth}}{
        \cellcolor[rgb]{0.9,0.9,0.9}{
            \makecell[{{p{\linewidth}}}]{
                \texttt{\tiny{[P1$\langle$GM]}}
                \texttt{Il tuo messaggio è stato inoltrato all'altro giocatore.} \\
\\ 
\\ 
\texttt{Stai attualmente giocando il round 6 di massimo 12.} \\
\texttt{Avete accumulato collettivamente 2 di 8 penalità.} \\
\texttt{L'altro giocatore ha inviato questo messaggio:} \\
\\ 
\\ 
\texttt{"Accordo di spostare 'L' a (4,3) per evitare conflitti. Procediamo con questa posizione. 😊"} \\
\\ 
\\ 
\texttt{Qual è il tuo prossimo comando?} \\
            }
        }
    }
    & \\ \\

    \theutterance \stepcounter{utterance}  
    & \multicolumn{4}{p{0.6\linewidth}}{
        \cellcolor[rgb]{0.9,0.9,0.9}{
            \makecell[{{p{\linewidth}}}]{
                \texttt{\tiny{[P1$\rangle$GM]}}
                \texttt{SPOSTA: L, (4,3)} \\
            }
        }
    }
    & & \\ \\

    \theutterance \stepcounter{utterance}  
    & & & \multicolumn{2}{p{0.3\linewidth}}{
        \cellcolor[rgb]{0.9,0.9,0.9}{
            \makecell[{{p{\linewidth}}}]{
                \texttt{\tiny{[GM$|$GM]}}
                \texttt{SPOSTA: L, (4,3)} \\
            }
        }
    }
    & & \\ \\

    \theutterance \stepcounter{utterance}  
    & & & \multicolumn{2}{p{0.3\linewidth}}{
        \cellcolor[rgb]{0.9,0.9,0.9}{
            \makecell[{{p{\linewidth}}}]{
                \texttt{\tiny{[GM$|$GM]}}
                \texttt{Penalità: (4,3) non è vuoto, ma contiene '│'.} \\
\\ 
\texttt{Avete accumulato collettivamente 3 di 8 penalità.} \\
\texttt{Riprova!} \\
            }
        }
    }
    & & \\ \\

    \theutterance \stepcounter{utterance}  
    & & \multicolumn{4}{p{0.6\linewidth}}{
        \cellcolor[rgb]{0.9,0.9,0.9}{
            \makecell[{{p{\linewidth}}}]{
                \texttt{\tiny{[P1$\langle$GM]}}
                \texttt{Penalità: (4,3) non è vuoto, ma contiene '│'.} \\
\\ 
\texttt{Avete accumulato collettivamente 3 di 8 penalità.} \\
\texttt{Riprova!} \\
            }
        }
    }
    & \\ \\

    \theutterance \stepcounter{utterance}  
    & \multicolumn{4}{p{0.6\linewidth}}{
        \cellcolor[rgb]{0.9,0.9,0.9}{
            \makecell[{{p{\linewidth}}}]{
                \texttt{\tiny{[P1$\rangle$GM]}}
                \texttt{COMUNICA: Scusa per l'errore! La posizione (4,3) è occupata. Proponiamo di spostare 'L' in una cella vicina e libera, come (3,3) o (5,3). Tuttavia, (3,3) è già occupata da 'C'. Quale posizione preferisci? 😊} \\
            }
        }
    }
    & & \\ \\

    \theutterance \stepcounter{utterance}  
    & & & \multicolumn{2}{p{0.3\linewidth}}{
        \cellcolor[rgb]{0.9,0.9,0.9}{
            \makecell[{{p{\linewidth}}}]{
                \texttt{\tiny{[GM$|$GM]}}
                \texttt{COMUNICA: Scusa per l'errore! La posizione (4,3) è occupata. Proponiamo di spostare 'L' in una cella vicina e libera, come (3,3) o (5,3). Tuttavia, (3,3) è già occupata da 'C'. Quale posizione preferisci? 😊} \\
            }
        }
    }
    & & \\ \\

    \theutterance \stepcounter{utterance}  
    & & \multicolumn{4}{p{0.6\linewidth}}{
        \cellcolor[rgb]{0.9,0.9,0.9}{
            \makecell[{{p{\linewidth}}}]{
                \texttt{\tiny{[GM$\rangle$P2]}}
                \texttt{Il tuo messaggio è stato inoltrato all'altro giocatore.} \\
\\ 
\\ 
\texttt{Stai attualmente giocando il round 7 di massimo 12.} \\
\texttt{Avete accumulato collettivamente 3 di 8 penalità.} \\
\texttt{L'altro giocatore ha inviato questo messaggio:} \\
\\ 
\\ 
\texttt{"Scusa per l'errore! La posizione (4,3) è occupata. Proponiamo di spostare 'L' in una cella vicina e libera, come (3,3) o (5,3). Tuttavia, (3,3) è già occupata da 'C'. Quale posizione preferisci? 😊"} \\
\\ 
\\ 
\texttt{Qual è il tuo prossimo comando?} \\
            }
        }
    }
    & \\ \\

    \theutterance \stepcounter{utterance}  
    & & & \multicolumn{4}{p{0.6\linewidth}}{
        \cellcolor[rgb]{0.9,0.9,0.9}{
            \makecell[{{p{\linewidth}}}]{
                \texttt{\tiny{[GM$\langle$P2]}}
                \texttt{COMUNICA: Accordo di spostare 'L' a (5,3). Procediamo con questa posizione. 😊} \\
            }
        }
    }
     \\ \\

    \theutterance \stepcounter{utterance}  
    & & & \multicolumn{2}{p{0.3\linewidth}}{
        \cellcolor[rgb]{0.9,0.9,0.9}{
            \makecell[{{p{\linewidth}}}]{
                \texttt{\tiny{[GM$|$GM]}}
                \texttt{COMUNICA: Accordo di spostare 'L' a (5,3). Procediamo con questa posizione. 😊} \\
            }
        }
    }
    & & \\ \\

    \theutterance \stepcounter{utterance}  
    & & \multicolumn{4}{p{0.6\linewidth}}{
        \cellcolor[rgb]{0.9,0.9,0.9}{
            \makecell[{{p{\linewidth}}}]{
                \texttt{\tiny{[P1$\langle$GM]}}
                \texttt{Il tuo messaggio è stato inoltrato all'altro giocatore.} \\
\\ 
\\ 
\texttt{Stai attualmente giocando il round 7 di massimo 12.} \\
\texttt{Avete accumulato collettivamente 3 di 8 penalità.} \\
\texttt{L'altro giocatore ha inviato questo messaggio:} \\
\\ 
\\ 
\texttt{"Accordo di spostare 'L' a (5,3). Procediamo con questa posizione. 😊"} \\
\\ 
\\ 
\texttt{Qual è il tuo prossimo comando?} \\
            }
        }
    }
    & \\ \\

    \theutterance \stepcounter{utterance}  
    & \multicolumn{4}{p{0.6\linewidth}}{
        \cellcolor[rgb]{0.9,0.9,0.9}{
            \makecell[{{p{\linewidth}}}]{
                \texttt{\tiny{[P1$\rangle$GM]}}
                \texttt{SPOSTA: L, (5,3)} \\
            }
        }
    }
    & & \\ \\

    \theutterance \stepcounter{utterance}  
    & & & \multicolumn{2}{p{0.3\linewidth}}{
        \cellcolor[rgb]{0.9,0.9,0.9}{
            \makecell[{{p{\linewidth}}}]{
                \texttt{\tiny{[GM$|$GM]}}
                \texttt{SPOSTA: L, (5,3)} \\
            }
        }
    }
    & & \\ \\

    \theutterance \stepcounter{utterance}  
    & & & \multicolumn{2}{p{0.3\linewidth}}{
        \cellcolor[rgb]{0.9,0.9,0.9}{
            \makecell[{{p{\linewidth}}}]{
                \texttt{\tiny{[GM$|$GM]}}
                \texttt{Spostato 'L' in (5,3) con successo. La tua griglia aggiornata appare così:} \\
\texttt{```} \\
\texttt{ 1234567} \\
\texttt{╔═╤═╤═╤═╗} \\
\texttt{║◌│◌│◌│◌║ 1} \\
\texttt{╟─┴─┼─┤◌║ 2} \\
\texttt{║◌◌C│L└─╢ 3} \\
\texttt{╟───┘◌◌◌║ 4} \\
\texttt{║◌◌◌◌◌┌─╢ 5} \\
\texttt{║◌◌◌◌◌│◌║ 6} \\
\texttt{║◌◌◌◌P│◌║ 7} \\
\texttt{╚═════╧═╝} \\
\texttt{```} \\
            }
        }
    }
    & & \\ \\

    \theutterance \stepcounter{utterance}  
    & & \multicolumn{4}{p{0.6\linewidth}}{
        \cellcolor[rgb]{0.9,0.9,0.9}{
            \makecell[{{p{\linewidth}}}]{
                \texttt{\tiny{[GM$\rangle$P2]}}
                \texttt{Il tuo messaggio è stato inoltrato all'altro giocatore.} \\
\\ 
\\ 
\texttt{Stai attualmente giocando il round 8 di massimo 12.} \\
\texttt{Avete accumulato collettivamente 3 di 8 penalità.} \\
\texttt{L'altro giocatore ha spostato un oggetto sulla sua griglia.} \\
\\ 
\\ 
\texttt{Qual è il tuo prossimo comando?} \\
            }
        }
    }
    & \\ \\

    \theutterance \stepcounter{utterance}  
    & & & \multicolumn{4}{p{0.6\linewidth}}{
        \cellcolor[rgb]{0.9,0.9,0.9}{
            \makecell[{{p{\linewidth}}}]{
                \texttt{\tiny{[GM$\langle$P2]}}
                \texttt{SPOSTA: L, (5,3)} \\
            }
        }
    }
     \\ \\

    \theutterance \stepcounter{utterance}  
    & & & \multicolumn{2}{p{0.3\linewidth}}{
        \cellcolor[rgb]{0.9,0.9,0.9}{
            \makecell[{{p{\linewidth}}}]{
                \texttt{\tiny{[GM$|$GM]}}
                \texttt{SPOSTA: L, (5,3)} \\
            }
        }
    }
    & & \\ \\

    \theutterance \stepcounter{utterance}  
    & & & \multicolumn{2}{p{0.3\linewidth}}{
        \cellcolor[rgb]{0.9,0.9,0.9}{
            \makecell[{{p{\linewidth}}}]{
                \texttt{\tiny{[GM$|$GM]}}
                \texttt{Spostato 'L' in (5,3) con successo. La tua griglia aggiornata appare così:} \\
\texttt{```} \\
\texttt{ 1234567} \\
\texttt{╔═╤═╤═╤═╗} \\
\texttt{║◌│◌│◌│◌║ 1} \\
\texttt{╟─┴─┼─┤◌║ 2} \\
\texttt{║◌◌C│L└─╢ 3} \\
\texttt{╟───┘◌◌◌║ 4} \\
\texttt{║P◌◌◌◌┌─╢ 5} \\
\texttt{║◌◌◌◌◌│◌║ 6} \\
\texttt{║◌◌◌◌◌│◌║ 7} \\
\texttt{╚═════╧═╝} \\
\texttt{```} \\
            }
        }
    }
    & & \\ \\

    \theutterance \stepcounter{utterance}  
    & & \multicolumn{4}{p{0.6\linewidth}}{
        \cellcolor[rgb]{0.9,0.9,0.9}{
            \makecell[{{p{\linewidth}}}]{
                \texttt{\tiny{[P1$\langle$GM]}}
                \texttt{Spostato 'L' in (5,3) con successo. La tua griglia aggiornata appare così:} \\
\texttt{```} \\
\texttt{ 1234567} \\
\texttt{╔═╤═╤═╤═╗} \\
\texttt{║◌│◌│◌│◌║ 1} \\
\texttt{╟─┴─┼─┤◌║ 2} \\
\texttt{║◌◌C│L└─╢ 3} \\
\texttt{╟───┘◌◌◌║ 4} \\
\texttt{║◌◌◌◌◌┌─╢ 5} \\
\texttt{║◌◌◌◌◌│◌║ 6} \\
\texttt{║◌◌◌◌P│◌║ 7} \\
\texttt{╚═════╧═╝} \\
\texttt{```} \\
\texttt{Stai attualmente giocando il round 8 di massimo 12.} \\
\texttt{Avete accumulato collettivamente 3 di 8 penalità.} \\
\texttt{L'altro giocatore ha spostato un oggetto sulla sua griglia.} \\
\\ 
\\ 
\texttt{Qual è il tuo prossimo comando?} \\
            }
        }
    }
    & \\ \\

    \theutterance \stepcounter{utterance}  
    & \multicolumn{4}{p{0.6\linewidth}}{
        \cellcolor[rgb]{0.9,0.9,0.9}{
            \makecell[{{p{\linewidth}}}]{
                \texttt{\tiny{[P1$\rangle$GM]}}
                \texttt{COMUNICA: Ottimo! Ora che 'C' e 'L' sono posizionati, possiamo concentrarci su 'P'. Proponiamo di spostare 'P' verso (4,4) o (5,4) per allinearlo con gli altri. Quale posizione preferisci? 😊} \\
            }
        }
    }
    & & \\ \\

    \theutterance \stepcounter{utterance}  
    & & & \multicolumn{2}{p{0.3\linewidth}}{
        \cellcolor[rgb]{0.9,0.9,0.9}{
            \makecell[{{p{\linewidth}}}]{
                \texttt{\tiny{[GM$|$GM]}}
                \texttt{COMUNICA: Ottimo! Ora che 'C' e 'L' sono posizionati, possiamo concentrarci su 'P'. Proponiamo di spostare 'P' verso (4,4) o (5,4) per allinearlo con gli altri. Quale posizione preferisci? 😊} \\
            }
        }
    }
    & & \\ \\

    \theutterance \stepcounter{utterance}  
    & & \multicolumn{4}{p{0.6\linewidth}}{
        \cellcolor[rgb]{0.9,0.9,0.9}{
            \makecell[{{p{\linewidth}}}]{
                \texttt{\tiny{[GM$\rangle$P2]}}
                \texttt{Spostato 'L' in (5,3) con successo. La tua griglia aggiornata appare così:} \\
\texttt{```} \\
\texttt{ 1234567} \\
\texttt{╔═╤═╤═╤═╗} \\
\texttt{║◌│◌│◌│◌║ 1} \\
\texttt{╟─┴─┼─┤◌║ 2} \\
\texttt{║◌◌C│L└─╢ 3} \\
\texttt{╟───┘◌◌◌║ 4} \\
\texttt{║P◌◌◌◌┌─╢ 5} \\
\texttt{║◌◌◌◌◌│◌║ 6} \\
\texttt{║◌◌◌◌◌│◌║ 7} \\
\texttt{╚═════╧═╝} \\
\texttt{```} \\
\texttt{Stai attualmente giocando il round 9 di massimo 12.} \\
\texttt{Avete accumulato collettivamente 3 di 8 penalità.} \\
\texttt{L'altro giocatore ha inviato questo messaggio:} \\
\\ 
\\ 
\texttt{"Ottimo! Ora che 'C' e 'L' sono posizionati, possiamo concentrarci su 'P'. Proponiamo di spostare 'P' verso (4,4) o (5,4) per allinearlo con gli altri. Quale posizione preferisci? 😊"} \\
\\ 
\\ 
\texttt{Qual è il tuo prossimo comando?} \\
            }
        }
    }
    & \\ \\

    \theutterance \stepcounter{utterance}  
    & & & \multicolumn{4}{p{0.6\linewidth}}{
        \cellcolor[rgb]{0.9,0.9,0.9}{
            \makecell[{{p{\linewidth}}}]{
                \texttt{\tiny{[GM$\langle$P2]}}
                \texttt{COMUNICA: Proponiamo di spostare 'P' a (4,4) per allinearlo con gli altri oggetti. Se la posizione è occupata, possiamo adattare il piano. 😊} \\
            }
        }
    }
     \\ \\

    \theutterance \stepcounter{utterance}  
    & & & \multicolumn{2}{p{0.3\linewidth}}{
        \cellcolor[rgb]{0.9,0.9,0.9}{
            \makecell[{{p{\linewidth}}}]{
                \texttt{\tiny{[GM$|$GM]}}
                \texttt{COMUNICA: Proponiamo di spostare 'P' a (4,4) per allinearlo con gli altri oggetti. Se la posizione è occupata, possiamo adattare il piano. 😊} \\
            }
        }
    }
    & & \\ \\

    \theutterance \stepcounter{utterance}  
    & & \multicolumn{4}{p{0.6\linewidth}}{
        \cellcolor[rgb]{0.9,0.9,0.9}{
            \makecell[{{p{\linewidth}}}]{
                \texttt{\tiny{[P1$\langle$GM]}}
                \texttt{Il tuo messaggio è stato inoltrato all'altro giocatore.} \\
\\ 
\\ 
\texttt{Stai attualmente giocando il round 9 di massimo 12.} \\
\texttt{Avete accumulato collettivamente 3 di 8 penalità.} \\
\texttt{L'altro giocatore ha inviato questo messaggio:} \\
\\ 
\\ 
\texttt{"Proponiamo di spostare 'P' a (4,4) per allinearlo con gli altri oggetti. Se la posizione è occupata, possiamo adattare il piano. 😊"} \\
\\ 
\\ 
\texttt{Qual è il tuo prossimo comando?} \\
            }
        }
    }
    & \\ \\

    \theutterance \stepcounter{utterance}  
    & \multicolumn{4}{p{0.6\linewidth}}{
        \cellcolor[rgb]{0.9,0.9,0.9}{
            \makecell[{{p{\linewidth}}}]{
                \texttt{\tiny{[P1$\rangle$GM]}}
                \texttt{SPOSTA: P, (4,4)} \\
            }
        }
    }
    & & \\ \\

    \theutterance \stepcounter{utterance}  
    & & & \multicolumn{2}{p{0.3\linewidth}}{
        \cellcolor[rgb]{0.9,0.9,0.9}{
            \makecell[{{p{\linewidth}}}]{
                \texttt{\tiny{[GM$|$GM]}}
                \texttt{SPOSTA: P, (4,4)} \\
            }
        }
    }
    & & \\ \\

    \theutterance \stepcounter{utterance}  
    & & & \multicolumn{2}{p{0.3\linewidth}}{
        \cellcolor[rgb]{0.9,0.9,0.9}{
            \makecell[{{p{\linewidth}}}]{
                \texttt{\tiny{[GM$|$GM]}}
                \texttt{Penalità: (4,4) non è vuoto, ma contiene '┘'.} \\
\\ 
\texttt{Avete accumulato collettivamente 4 di 8 penalità.} \\
\texttt{Riprova!} \\
            }
        }
    }
    & & \\ \\

    \theutterance \stepcounter{utterance}  
    & & \multicolumn{4}{p{0.6\linewidth}}{
        \cellcolor[rgb]{0.9,0.9,0.9}{
            \makecell[{{p{\linewidth}}}]{
                \texttt{\tiny{[P1$\langle$GM]}}
                \texttt{Penalità: (4,4) non è vuoto, ma contiene '┘'.} \\
\\ 
\texttt{Avete accumulato collettivamente 4 di 8 penalità.} \\
\texttt{Riprova!} \\
            }
        }
    }
    & \\ \\

    \theutterance \stepcounter{utterance}  
    & \multicolumn{4}{p{0.6\linewidth}}{
        \cellcolor[rgb]{0.9,0.9,0.9}{
            \makecell[{{p{\linewidth}}}]{
                \texttt{\tiny{[P1$\rangle$GM]}}
                \texttt{COMUNICA: Scusa per l'errore! La posizione (4,4) è ancora occupata. Proponiamo di spostare 'P' in una cella vicina e libera, come (5,4) o (4,5). Quale posizione preferisci? 😊} \\
            }
        }
    }
    & & \\ \\

    \theutterance \stepcounter{utterance}  
    & & & \multicolumn{2}{p{0.3\linewidth}}{
        \cellcolor[rgb]{0.9,0.9,0.9}{
            \makecell[{{p{\linewidth}}}]{
                \texttt{\tiny{[GM$|$GM]}}
                \texttt{COMUNICA: Scusa per l'errore! La posizione (4,4) è ancora occupata. Proponiamo di spostare 'P' in una cella vicina e libera, come (5,4) o (4,5). Quale posizione preferisci? 😊} \\
            }
        }
    }
    & & \\ \\

    \theutterance \stepcounter{utterance}  
    & & \multicolumn{4}{p{0.6\linewidth}}{
        \cellcolor[rgb]{0.9,0.9,0.9}{
            \makecell[{{p{\linewidth}}}]{
                \texttt{\tiny{[GM$\rangle$P2]}}
                \texttt{Il tuo messaggio è stato inoltrato all'altro giocatore.} \\
\\ 
\\ 
\texttt{Stai attualmente giocando il round 10 di massimo 12.} \\
\texttt{Avete accumulato collettivamente 4 di 8 penalità.} \\
\texttt{L'altro giocatore ha inviato questo messaggio:} \\
\\ 
\\ 
\texttt{"Scusa per l'errore! La posizione (4,4) è ancora occupata. Proponiamo di spostare 'P' in una cella vicina e libera, come (5,4) o (4,5). Quale posizione preferisci? 😊"} \\
\\ 
\\ 
\texttt{Qual è il tuo prossimo comando?} \\
            }
        }
    }
    & \\ \\

    \theutterance \stepcounter{utterance}  
    & & & \multicolumn{4}{p{0.6\linewidth}}{
        \cellcolor[rgb]{0.9,0.9,0.9}{
            \makecell[{{p{\linewidth}}}]{
                \texttt{\tiny{[GM$\langle$P2]}}
                \texttt{COMUNICA: Accordo di spostare 'P' a (5,4) per evitare conflitti. Procediamo con questa posizione. 😊} \\
            }
        }
    }
     \\ \\

    \theutterance \stepcounter{utterance}  
    & & & \multicolumn{2}{p{0.3\linewidth}}{
        \cellcolor[rgb]{0.9,0.9,0.9}{
            \makecell[{{p{\linewidth}}}]{
                \texttt{\tiny{[GM$|$GM]}}
                \texttt{COMUNICA: Accordo di spostare 'P' a (5,4) per evitare conflitti. Procediamo con questa posizione. 😊} \\
            }
        }
    }
    & & \\ \\

    \theutterance \stepcounter{utterance}  
    & & \multicolumn{4}{p{0.6\linewidth}}{
        \cellcolor[rgb]{0.9,0.9,0.9}{
            \makecell[{{p{\linewidth}}}]{
                \texttt{\tiny{[P1$\langle$GM]}}
                \texttt{Il tuo messaggio è stato inoltrato all'altro giocatore.} \\
\\ 
\\ 
\texttt{Stai attualmente giocando il round 10 di massimo 12.} \\
\texttt{Avete accumulato collettivamente 4 di 8 penalità.} \\
\texttt{L'altro giocatore ha inviato questo messaggio:} \\
\\ 
\\ 
\texttt{"Accordo di spostare 'P' a (5,4) per evitare conflitti. Procediamo con questa posizione. 😊"} \\
\\ 
\\ 
\texttt{Qual è il tuo prossimo comando?} \\
            }
        }
    }
    & \\ \\

    \theutterance \stepcounter{utterance}  
    & \multicolumn{4}{p{0.6\linewidth}}{
        \cellcolor[rgb]{0.9,0.9,0.9}{
            \makecell[{{p{\linewidth}}}]{
                \texttt{\tiny{[P1$\rangle$GM]}}
                \texttt{SPOSTA: P, (5,4)} \\
            }
        }
    }
    & & \\ \\

    \theutterance \stepcounter{utterance}  
    & & & \multicolumn{2}{p{0.3\linewidth}}{
        \cellcolor[rgb]{0.9,0.9,0.9}{
            \makecell[{{p{\linewidth}}}]{
                \texttt{\tiny{[GM$|$GM]}}
                \texttt{SPOSTA: P, (5,4)} \\
            }
        }
    }
    & & \\ \\

    \theutterance \stepcounter{utterance}  
    & & & \multicolumn{2}{p{0.3\linewidth}}{
        \cellcolor[rgb]{0.9,0.9,0.9}{
            \makecell[{{p{\linewidth}}}]{
                \texttt{\tiny{[GM$|$GM]}}
                \texttt{Spostato 'P' in (5,4) con successo. La tua griglia aggiornata appare così:} \\
\texttt{```} \\
\texttt{ 1234567} \\
\texttt{╔═╤═╤═╤═╗} \\
\texttt{║◌│◌│◌│◌║ 1} \\
\texttt{╟─┴─┼─┤◌║ 2} \\
\texttt{║◌◌C│L└─╢ 3} \\
\texttt{╟───┘P◌◌║ 4} \\
\texttt{║◌◌◌◌◌┌─╢ 5} \\
\texttt{║◌◌◌◌◌│◌║ 6} \\
\texttt{║◌◌◌◌◌│◌║ 7} \\
\texttt{╚═════╧═╝} \\
\texttt{```} \\
            }
        }
    }
    & & \\ \\

    \theutterance \stepcounter{utterance}  
    & & \multicolumn{4}{p{0.6\linewidth}}{
        \cellcolor[rgb]{0.9,0.9,0.9}{
            \makecell[{{p{\linewidth}}}]{
                \texttt{\tiny{[GM$\rangle$P2]}}
                \texttt{Il tuo messaggio è stato inoltrato all'altro giocatore.} \\
\\ 
\\ 
\texttt{Stai attualmente giocando il round 11 di massimo 12.} \\
\texttt{Avete accumulato collettivamente 4 di 8 penalità.} \\
\texttt{L'altro giocatore ha spostato un oggetto sulla sua griglia.} \\
\\ 
\\ 
\texttt{Qual è il tuo prossimo comando?} \\
            }
        }
    }
    & \\ \\

    \theutterance \stepcounter{utterance}  
    & & & \multicolumn{4}{p{0.6\linewidth}}{
        \cellcolor[rgb]{0.9,0.9,0.9}{
            \makecell[{{p{\linewidth}}}]{
                \texttt{\tiny{[GM$\langle$P2]}}
                \texttt{SPOSTA: P, (5,4)} \\
            }
        }
    }
     \\ \\

    \theutterance \stepcounter{utterance}  
    & & & \multicolumn{2}{p{0.3\linewidth}}{
        \cellcolor[rgb]{0.9,0.9,0.9}{
            \makecell[{{p{\linewidth}}}]{
                \texttt{\tiny{[GM$|$GM]}}
                \texttt{SPOSTA: P, (5,4)} \\
            }
        }
    }
    & & \\ \\

    \theutterance \stepcounter{utterance}  
    & & & \multicolumn{2}{p{0.3\linewidth}}{
        \cellcolor[rgb]{0.9,0.9,0.9}{
            \makecell[{{p{\linewidth}}}]{
                \texttt{\tiny{[GM$|$GM]}}
                \texttt{Spostato 'P' in (5,4) con successo. La tua griglia aggiornata appare così:} \\
\texttt{```} \\
\texttt{ 1234567} \\
\texttt{╔═╤═╤═╤═╗} \\
\texttt{║◌│◌│◌│◌║ 1} \\
\texttt{╟─┴─┼─┤◌║ 2} \\
\texttt{║◌◌C│L└─╢ 3} \\
\texttt{╟───┘P◌◌║ 4} \\
\texttt{║◌◌◌◌◌┌─╢ 5} \\
\texttt{║◌◌◌◌◌│◌║ 6} \\
\texttt{║◌◌◌◌◌│◌║ 7} \\
\texttt{╚═════╧═╝} \\
\texttt{```} \\
            }
        }
    }
    & & \\ \\

    \theutterance \stepcounter{utterance}  
    & & \multicolumn{4}{p{0.6\linewidth}}{
        \cellcolor[rgb]{0.9,0.9,0.9}{
            \makecell[{{p{\linewidth}}}]{
                \texttt{\tiny{[P1$\langle$GM]}}
                \texttt{Spostato 'P' in (5,4) con successo. La tua griglia aggiornata appare così:} \\
\texttt{```} \\
\texttt{ 1234567} \\
\texttt{╔═╤═╤═╤═╗} \\
\texttt{║◌│◌│◌│◌║ 1} \\
\texttt{╟─┴─┼─┤◌║ 2} \\
\texttt{║◌◌C│L└─╢ 3} \\
\texttt{╟───┘P◌◌║ 4} \\
\texttt{║◌◌◌◌◌┌─╢ 5} \\
\texttt{║◌◌◌◌◌│◌║ 6} \\
\texttt{║◌◌◌◌◌│◌║ 7} \\
\texttt{╚═════╧═╝} \\
\texttt{```} \\
\texttt{Stai attualmente giocando il round 11 di massimo 12.} \\
\texttt{Avete accumulato collettivamente 4 di 8 penalità.} \\
\texttt{L'altro giocatore ha spostato un oggetto sulla sua griglia.} \\
\\ 
\\ 
\texttt{Qual è il tuo prossimo comando?} \\
            }
        }
    }
    & \\ \\

    \theutterance \stepcounter{utterance}  
    & \multicolumn{4}{p{0.6\linewidth}}{
        \cellcolor[rgb]{0.9,0.9,0.9}{
            \makecell[{{p{\linewidth}}}]{
                \texttt{\tiny{[P1$\rangle$GM]}}
                \texttt{COMUNICA: Ottimo! Ora tutti gli oggetti (C, L, P) sono posizionati in coordinate vicine: C a (3,3), L a (5,3) e P a (5,4). Proponiamo di finalizzare il gioco, poiché le posizioni sono coerenti e vicine. Possiamo inviare COMUNICA: finito? per confermare. 😊} \\
            }
        }
    }
    & & \\ \\

    \theutterance \stepcounter{utterance}  
    & & & \multicolumn{2}{p{0.3\linewidth}}{
        \cellcolor[rgb]{0.9,0.9,0.9}{
            \makecell[{{p{\linewidth}}}]{
                \texttt{\tiny{[GM$|$GM]}}
                \texttt{COMUNICA: Ottimo! Ora tutti gli oggetti (C, L, P) sono posizionati in coordinate vicine: C a (3,3), L a (5,3) e P a (5,4). Proponiamo di finalizzare il gioco, poiché le posizioni sono coerenti e vicine. Possiamo inviare COMUNICA: finito? per confermare. 😊} \\
            }
        }
    }
    & & \\ \\

    \theutterance \stepcounter{utterance}  
    & & & \multicolumn{2}{p{0.3\linewidth}}{
        \cellcolor[rgb]{0.9,0.9,0.9}{
            \makecell[{{p{\linewidth}}}]{
                \texttt{\tiny{[GM$|$GM]}}
                \texttt{Invalid format: head is not empty: 'COMUNICA: Ottimo! Ora tutti gli oggetti (C, L, P) sono posizionati in coordinate vicine: C a (3,3), L a (5,3) e P a (5,4). Proponiamo di finalizzare il gioco, poiché le posizioni sono coerenti e vicine. Possiamo inviare '} \\
            }
        }
    }
    & & \\ \\

    \theutterance \stepcounter{utterance}  
    & & \multicolumn{4}{p{0.6\linewidth}}{
        \cellcolor[rgb]{0.9,0.9,0.9}{
            \makecell[{{p{\linewidth}}}]{
                \texttt{\tiny{[P1$\langle$GM]}}
                \texttt{Penalità: Il tuo messaggio non deve contenere nulla prima del comando!} \\
\\ 
\texttt{Assicurati che la tua risposta contenga solo `COMUNICA: <MESSAGE>` o `SPOSTA: <OBJECT>, (<X>, <Y>)`, e nient'altro!} \\
\\ 
\texttt{Avete accumulato collettivamente 5 di 8 penalità.} \\
\texttt{Riprova!} \\
            }
        }
    }
    & \\ \\

    \theutterance \stepcounter{utterance}  
    & \multicolumn{4}{p{0.6\linewidth}}{
        \cellcolor[rgb]{0.9,0.9,0.9}{
            \makecell[{{p{\linewidth}}}]{
                \texttt{\tiny{[P1$\rangle$GM]}}
                \texttt{COMUNICA: finito?} \\
            }
        }
    }
    & & \\ \\

    \theutterance \stepcounter{utterance}  
    & & & \multicolumn{2}{p{0.3\linewidth}}{
        \cellcolor[rgb]{0.9,0.9,0.9}{
            \makecell[{{p{\linewidth}}}]{
                \texttt{\tiny{[GM$|$GM]}}
                \texttt{COMUNICA: finito?} \\
            }
        }
    }
    & & \\ \\

    \theutterance \stepcounter{utterance}  
    & & \multicolumn{4}{p{0.6\linewidth}}{
        \cellcolor[rgb]{0.9,0.9,0.9}{
            \makecell[{{p{\linewidth}}}]{
                \texttt{\tiny{[GM$\rangle$P2]}}
                \texttt{Spostato 'P' in (5,4) con successo. La tua griglia aggiornata appare così:} \\
\texttt{```} \\
\texttt{ 1234567} \\
\texttt{╔═╤═╤═╤═╗} \\
\texttt{║◌│◌│◌│◌║ 1} \\
\texttt{╟─┴─┼─┤◌║ 2} \\
\texttt{║◌◌C│L└─╢ 3} \\
\texttt{╟───┘P◌◌║ 4} \\
\texttt{║◌◌◌◌◌┌─╢ 5} \\
\texttt{║◌◌◌◌◌│◌║ 6} \\
\texttt{║◌◌◌◌◌│◌║ 7} \\
\texttt{╚═════╧═╝} \\
\texttt{```} \\
\texttt{Stai attualmente giocando il round 12 di massimo 12.} \\
\texttt{Avete accumulato collettivamente 5 di 8 penalità.} \\
\texttt{L'altro giocatore ha inviato questo messaggio:} \\
\\ 
\\ 
\texttt{"finito?"} \\
\\ 
\\ 
\texttt{Qual è il tuo prossimo comando?} \\
            }
        }
    }
    & \\ \\

    \theutterance \stepcounter{utterance}  
    & & & \multicolumn{4}{p{0.6\linewidth}}{
        \cellcolor[rgb]{0.9,0.9,0.9}{
            \makecell[{{p{\linewidth}}}]{
                \texttt{\tiny{[GM$\langle$P2]}}
                \texttt{COMUNICA: finito!} \\
            }
        }
    }
     \\ \\

    \theutterance \stepcounter{utterance}  
    & & & \multicolumn{2}{p{0.3\linewidth}}{
        \cellcolor[rgb]{0.9,0.9,0.9}{
            \makecell[{{p{\linewidth}}}]{
                \texttt{\tiny{[GM$|$GM]}}
                \texttt{COMUNICA: finito!} \\
            }
        }
    }
    & & \\ \\

    \theutterance \stepcounter{utterance}  
    & & & \multicolumn{2}{p{0.3\linewidth}}{
        \cellcolor[rgb]{0.9,0.9,0.9}{
            \makecell[{{p{\linewidth}}}]{
                \texttt{\tiny{[GM$|$GM]}}
                \texttt{true} \\
            }
        }
    }
    & & \\ \\

    \theutterance \stepcounter{utterance}  
    & & & \multicolumn{2}{p{0.3\linewidth}}{
        \cellcolor[rgb]{0.9,0.9,0.9}{
            \makecell[{{p{\linewidth}}}]{
                \texttt{\tiny{[GM$|$GM]}}
                \texttt{* success: True} \\
\texttt{* lose: False} \\
\texttt{* aborted: False} \\
\texttt{{-}{-}{-}{-}{-}{-}{-}} \\
\texttt{* Shifts: 2.00} \\
\texttt{* Max Shifts: 4.00} \\
\texttt{* Min Shifts: 2.00} \\
\texttt{* End Distance Sum: 0.00} \\
\texttt{* Init Distance Sum: 9.47} \\
\texttt{* Expected Distance Sum: 12.57} \\
\texttt{* Penalties: 5.00} \\
\texttt{* Max Penalties: 8.00} \\
\texttt{* Rounds: 12.00} \\
\texttt{* Max Rounds: 12.00} \\
\texttt{* Object Count: 3.00} \\
            }
        }
    }
    & & \\ \\

    \theutterance \stepcounter{utterance}  
    & & & \multicolumn{2}{p{0.3\linewidth}}{
        \cellcolor[rgb]{0.9,0.9,0.9}{
            \makecell[{{p{\linewidth}}}]{
                \texttt{\tiny{[GM$|$GM]}}
                \texttt{Initial Grids:} \\
\texttt{```} \\
\texttt{Player 1:     Player 2:} \\
\texttt{ 1234567       1234567} \\
\texttt{╔═╤═╤═╤═╗     ╔═╤═╤═╤═╗} \\
\texttt{║◌│◌│◌│◌║ 1   ║◌│◌│L│◌║ 1} \\
\texttt{╟─┴─┼─┤◌║ 2   ╟─┴─┼─┤◌║ 2} \\
\texttt{║C◌◌│◌└─╢ 3   ║C◌◌│◌└─╢ 3} \\
\texttt{╟───┘◌◌◌║ 4   ╟───┘◌◌◌║ 4} \\
\texttt{║◌L◌◌◌┌─╢ 5   ║P◌◌◌◌┌─╢ 5} \\
\texttt{║◌◌◌◌◌│◌║ 6   ║◌◌◌◌◌│◌║ 6} \\
\texttt{║◌◌◌◌P│◌║ 7   ║◌◌◌◌◌│◌║ 7} \\
\texttt{╚═════╧═╝     ╚═════╧═╝} \\
\texttt{```} \\
\\ 
\texttt{Final Grids:} \\
\texttt{```} \\
\texttt{Player 1:     Player 2:} \\
\texttt{ 1234567       1234567} \\
\texttt{╔═╤═╤═╤═╗     ╔═╤═╤═╤═╗} \\
\texttt{║◌│◌│◌│◌║ 1   ║◌│◌│◌│◌║ 1} \\
\texttt{╟─┴─┼─┤◌║ 2   ╟─┴─┼─┤◌║ 2} \\
\texttt{║◌◌C│L└─╢ 3   ║◌◌C│L└─╢ 3} \\
\texttt{╟───┘P◌◌║ 4   ╟───┘P◌◌║ 4} \\
\texttt{║◌◌◌◌◌┌─╢ 5   ║◌◌◌◌◌┌─╢ 5} \\
\texttt{║◌◌◌◌◌│◌║ 6   ║◌◌◌◌◌│◌║ 6} \\
\texttt{║◌◌◌◌◌│◌║ 7   ║◌◌◌◌◌│◌║ 7} \\
\texttt{╚═════╧═╝     ╚═════╧═╝} \\
\texttt{```} \\
            }
        }
    }
    & & \\ \\

\end{supertabular}
}

\end{document}
